\documentclass[11pt,a4paper]{article}
\usepackage[utf8]{inputenc}
\usepackage[spanish]{babel}
\usepackage{amsmath}
\usepackage{amssymb}
\usepackage{graphicx}
\usepackage{geometry}
\usepackage{caption}
\usepackage{subcaption}
\usepackage{booktabs}
\usepackage{tabularx}
\usepackage{hyperref}
\usepackage{float}
\addto\captionsspanish{\renewcommand{\tablename}{Tabla}}
\addto\captionsspanish{\renewcommand{\figurename}{Figura}}
\geometry{top=2cm, bottom=2cm, left=2.2cm, right=2.2cm}

\title{\textbf{Reporte Técnico: Análisis de SNR y Calidad de Señal (2026-09-23)}}
\author{Laboratorio de Sistemas Dinámicos (LSD) - Sistema Ñandú}
\date{2026-09-23}

\begin{document}
\maketitle

\tableofcontents
\newpage
\section{Condiciones Experimentales}
\begin{itemize}
    \item \textbf{Fecha:} 2026-09-23
    \item \textbf{Sujeto:} Sujeto1
    \item \textbf{Frecuencia de Muestreo:} 2000 Hz
    \item \textbf{Total de Registros Analizados:} 19
\end{itemize}

\section{Tabla de Relación Señal-Ruido (SNR) de Todas las Mediciones}
A continuación se presentan los valores de Relación Señal-Ruido promedio por pulso con su correspondiente dispersión ($	ext{SNR} \pm \sigma$) obtenidos para cada una de las pruebas registradas en la sesión:
\begin{table}[H]
\centering
\small
\setlength{\tabcolsep}{4.5pt}
\begin{tabular}{c l c c c c c}
\toprule
\textbf{Nº} & \textbf{Identificador} & \textbf{Vocal} & \textbf{Serie} & \multicolumn{3}{c}{\textbf{Relación Señal-Ruido ($\text{SNR} \pm \sigma$)}} \\
\cmidrule(lr){5-7}
& & & & \textbf{Ch0: masetero} & \textbf{Ch1: -} & \textbf{Ch2: orbi} \\
\midrule
1 & O\_Pruebat1\_Sujeto1 & O & Pruebat1 & $1.9 \pm 0.0$ & $20.7 \pm 9.6$ & $18.8 \pm 8.7$ \\
2 & A\_Pruebat1\_Sujeto1 & A & Pruebat1 & $2.2 \pm 0.2$ & $3.8 \pm 1.6$ & $3.8 \pm 1.5$ \\
3 & U\_Pruebat1\_Sujeto1 & U & Pruebat1 & $2.1 \pm 0.0$ & $12.5 \pm 6.5$ & $12.2 \pm 6.2$ \\
4 & E\_Pruebat1\_Sujeto1 & E & Pruebat1 & $1.9 \pm 0.0$ & $8.7 \pm 2.1$ & $8.1 \pm 2.0$ \\
5 & I\_Pruebat1\_Sujeto1 & I & Pruebat1 & $2.1 \pm 0.2$ & $9.5 \pm 4.5$ & $8.9 \pm 4.2$ \\
6 & E\_Prueba2\_Sujeto1 & E & Prueba2 & $2.0 \pm 0.1$ & $6.2 \pm 2.0$ & $5.4 \pm 1.6$ \\
7 & U\_Prueba2\_Sujeto1 & U & Prueba2 & $1.6 \pm 0.0$ & $16.0 \pm 7.3$ & $14.4 \pm 6.4$ \\
8 & I\_Prueba2\_Sujeto1 & I & Prueba2 & $1.5 \pm 0.2$ & $11.2 \pm 4.9$ & $9.4 \pm 4.0$ \\
9 & A\_Prueba2\_Sujeto1 & A & Prueba2 & $2.2 \pm 0.3$ & $5.2 \pm 1.6$ & $4.5 \pm 1.2$ \\
10 & O\_Prueba2\_Sujeto1 & O & Prueba2 & $1.8 \pm 0.0$ & $18.0 \pm 6.2$ & $16.0 \pm 5.6$ \\
11 & A\_Prueba3\_Sujeto1 & A & Prueba3 & $2.5 \pm 0.2$ & $4.9 \pm 1.1$ & $4.4 \pm 0.9$ \\
12 & I\_Prueba3\_Sujeto1 & I & Prueba3 & $2.0 \pm 0.2$ & $8.9 \pm 2.6$ & $8.0 \pm 2.3$ \\
13 & O\_Prueba3\_Sujeto1 & O & Prueba3 & $1.7 \pm 0.0$ & $12.8 \pm 3.9$ & $12.4 \pm 3.8$ \\
14 & U\_Prueba3\_Sujeto1 & U & Prueba3 & $1.4 \pm 0.0$ & $16.1 \pm 3.3$ & $14.4 \pm 2.9$ \\
15 & E\_Prueba3\_Sujeto1 & E & Prueba3 & $1.9 \pm 0.1$ & $5.5 \pm 1.5$ & $4.9 \pm 1.2$ \\
16 & O\_Prueba4\_Sujeto1 & O & Prueba4 & $1.9 \pm 0.1$ & $16.1 \pm 6.5$ & $26.0 \pm 9.8$ \\
17 & U\_Prueba4\_Sujeto1 & U & Prueba4 & $1.7 \pm 0.1$ & $14.4 \pm 6.4$ & $30.6 \pm 14.1$ \\
18 & O\_Prueba5\_Sujeto1 & O & Prueba5 & $1.7 \pm 0.0$ & $16.4 \pm 7.4$ & $34.0 \pm 16.1$ \\
19 & U\_Prueba5\_Sujeto1 & U & Prueba5 & $1.6 \pm 0.1$ & $20.0 \pm 9.7$ & $29.1 \pm 13.9$ \\
\bottomrule
\end{tabular}
\caption{Relación Señal-Ruido ($\text{SNR} \pm \sigma$) calculada por pulso para la totalidad de registros adquiridos en la sesión.}
\end{table}

\section{Resumen Estadístico de Relación Señal-Ruido (SNR)}
\begin{table}[H]
\centering
\begin{tabular}{lccc}
\toprule
\textbf{Vocal} & \textbf{Canal 0: masetero} & \textbf{Canal 1: -} & \textbf{Canal 2: orbi} \\
\midrule
Vocal A & $2.4 \pm 0.3$ & $4.7 \pm 1.5$ & $4.3 \pm 1.2$ \\
Vocal E & $1.9 \pm 0.1$ & $6.4 \pm 2.2$ & $5.7 \pm 2.0$ \\
Vocal I & $1.9 \pm 0.3$ & $9.5 \pm 3.8$ & $8.5 \pm 3.3$ \\
Vocal O & $1.8 \pm 0.1$ & $15.9 \pm 7.0$ & $19.5 \pm 11.8$ \\
Vocal U & $1.6 \pm 0.2$ & $15.9 \pm 6.7$ & $18.9 \pm 11.5$ \\
\bottomrule
\end{tabular}
\caption{Relación Señal-Ruido promedio ($\text{SNR} \pm \sigma$) registrada por vocal y canal.}
\end{table}

\begin{figure}[H]
\centering
\includegraphics[width=0.88\textwidth]{/home/santiago/repositorios/Nandu_SistemadeAdqusicionEMG/reportes_experimentos/evolucion_sesion/Sesion_SNR_Timeline.png}
\caption{Evolución temporal del SNR a lo largo de la sesión.}
\end{figure}

\newpage
\section{Desglose de Calidad de Señal, Recortes y Ruido Interpulso}
\subsection{Pruebat1}
\subsubsection{Vocal O: O\_Pruebat1\_Sujeto1}
\begin{figure}[H]
\centering
\includegraphics[width=0.88\textwidth]{/home/santiago/repositorios/Nandu_SistemadeAdqusicionEMG/EMG_desarrollo/base_de_datos_electrodos/2026-09-23/O_Pruebat1_Sujeto1/plot_paper_combined.png}
\caption{Análisis multimodal de la vocal O (O\_Pruebat1\_Sujeto1): Espectrograma de audio con pre-énfasis, oscilograma acústico, activación EMG (Supremo Tricanal) y espectrograma muscular RGB.}
\end{figure}

\subsubsection{Vocal A: A\_Pruebat1\_Sujeto1}
\begin{figure}[H]
\centering
\includegraphics[width=0.88\textwidth]{/home/santiago/repositorios/Nandu_SistemadeAdqusicionEMG/EMG_desarrollo/base_de_datos_electrodos/2026-09-23/A_Pruebat1_Sujeto1/plot_paper_combined.png}
\caption{Análisis multimodal de la vocal A (A\_Pruebat1\_Sujeto1): Espectrograma de audio con pre-énfasis, oscilograma acústico, activación EMG (Supremo Tricanal) y espectrograma muscular RGB.}
\end{figure}

\subsubsection{Vocal U: U\_Pruebat1\_Sujeto1}
\begin{figure}[H]
\centering
\includegraphics[width=0.88\textwidth]{/home/santiago/repositorios/Nandu_SistemadeAdqusicionEMG/EMG_desarrollo/base_de_datos_electrodos/2026-09-23/U_Pruebat1_Sujeto1/plot_paper_combined.png}
\caption{Análisis multimodal de la vocal U (U\_Pruebat1\_Sujeto1): Espectrograma de audio con pre-énfasis, oscilograma acústico, activación EMG (Supremo Tricanal) y espectrograma muscular RGB.}
\end{figure}

\subsubsection{Vocal E: E\_Pruebat1\_Sujeto1}
\begin{figure}[H]
\centering
\includegraphics[width=0.88\textwidth]{/home/santiago/repositorios/Nandu_SistemadeAdqusicionEMG/EMG_desarrollo/base_de_datos_electrodos/2026-09-23/E_Pruebat1_Sujeto1/plot_paper_combined.png}
\caption{Análisis multimodal de la vocal E (E\_Pruebat1\_Sujeto1): Espectrograma de audio con pre-énfasis, oscilograma acústico, activación EMG (Supremo Tricanal) y espectrograma muscular RGB.}
\end{figure}

\subsubsection{Vocal I: I\_Pruebat1\_Sujeto1}
\begin{figure}[H]
\centering
\includegraphics[width=0.88\textwidth]{/home/santiago/repositorios/Nandu_SistemadeAdqusicionEMG/EMG_desarrollo/base_de_datos_electrodos/2026-09-23/I_Pruebat1_Sujeto1/plot_paper_combined.png}
\caption{Análisis multimodal de la vocal I (I\_Pruebat1\_Sujeto1): Espectrograma de audio con pre-énfasis, oscilograma acústico, activación EMG (Supremo Tricanal) y espectrograma muscular RGB.}
\end{figure}

\subsection{Prueba2}
\subsubsection{Vocal E: E\_Prueba2\_Sujeto1}
\begin{figure}[H]
\centering
\includegraphics[width=0.88\textwidth]{/home/santiago/repositorios/Nandu_SistemadeAdqusicionEMG/EMG_desarrollo/base_de_datos_electrodos/2026-09-23/E_Prueba2_Sujeto1/plot_paper_combined.png}
\caption{Análisis multimodal de la vocal E (E\_Prueba2\_Sujeto1): Espectrograma de audio con pre-énfasis, oscilograma acústico, activación EMG (Supremo Tricanal) y espectrograma muscular RGB.}
\end{figure}

\subsubsection{Vocal U: U\_Prueba2\_Sujeto1}
\begin{figure}[H]
\centering
\includegraphics[width=0.88\textwidth]{/home/santiago/repositorios/Nandu_SistemadeAdqusicionEMG/EMG_desarrollo/base_de_datos_electrodos/2026-09-23/U_Prueba2_Sujeto1/plot_paper_combined.png}
\caption{Análisis multimodal de la vocal U (U\_Prueba2\_Sujeto1): Espectrograma de audio con pre-énfasis, oscilograma acústico, activación EMG (Supremo Tricanal) y espectrograma muscular RGB.}
\end{figure}

\subsubsection{Vocal I: I\_Prueba2\_Sujeto1}
\begin{figure}[H]
\centering
\includegraphics[width=0.88\textwidth]{/home/santiago/repositorios/Nandu_SistemadeAdqusicionEMG/EMG_desarrollo/base_de_datos_electrodos/2026-09-23/I_Prueba2_Sujeto1/plot_paper_combined.png}
\caption{Análisis multimodal de la vocal I (I\_Prueba2\_Sujeto1): Espectrograma de audio con pre-énfasis, oscilograma acústico, activación EMG (Supremo Tricanal) y espectrograma muscular RGB.}
\end{figure}

\subsubsection{Vocal A: A\_Prueba2\_Sujeto1}
\begin{figure}[H]
\centering
\includegraphics[width=0.88\textwidth]{/home/santiago/repositorios/Nandu_SistemadeAdqusicionEMG/EMG_desarrollo/base_de_datos_electrodos/2026-09-23/A_Prueba2_Sujeto1/plot_paper_combined.png}
\caption{Análisis multimodal de la vocal A (A\_Prueba2\_Sujeto1): Espectrograma de audio con pre-énfasis, oscilograma acústico, activación EMG (Supremo Tricanal) y espectrograma muscular RGB.}
\end{figure}

\subsubsection{Vocal O: O\_Prueba2\_Sujeto1}
\begin{figure}[H]
\centering
\includegraphics[width=0.88\textwidth]{/home/santiago/repositorios/Nandu_SistemadeAdqusicionEMG/EMG_desarrollo/base_de_datos_electrodos/2026-09-23/O_Prueba2_Sujeto1/plot_paper_combined.png}
\caption{Análisis multimodal de la vocal O (O\_Prueba2\_Sujeto1): Espectrograma de audio con pre-énfasis, oscilograma acústico, activación EMG (Supremo Tricanal) y espectrograma muscular RGB.}
\end{figure}

\subsection{Prueba3}
\subsubsection{Vocal A: A\_Prueba3\_Sujeto1}
\begin{figure}[H]
\centering
\includegraphics[width=0.88\textwidth]{/home/santiago/repositorios/Nandu_SistemadeAdqusicionEMG/EMG_desarrollo/base_de_datos_electrodos/2026-09-23/A_Prueba3_Sujeto1/plot_paper_combined.png}
\caption{Análisis multimodal de la vocal A (A\_Prueba3\_Sujeto1): Espectrograma de audio con pre-énfasis, oscilograma acústico, activación EMG (Supremo Tricanal) y espectrograma muscular RGB.}
\end{figure}

\subsubsection{Vocal I: I\_Prueba3\_Sujeto1}
\begin{figure}[H]
\centering
\includegraphics[width=0.88\textwidth]{/home/santiago/repositorios/Nandu_SistemadeAdqusicionEMG/EMG_desarrollo/base_de_datos_electrodos/2026-09-23/I_Prueba3_Sujeto1/plot_paper_combined.png}
\caption{Análisis multimodal de la vocal I (I\_Prueba3\_Sujeto1): Espectrograma de audio con pre-énfasis, oscilograma acústico, activación EMG (Supremo Tricanal) y espectrograma muscular RGB.}
\end{figure}

\subsubsection{Vocal O: O\_Prueba3\_Sujeto1}
\begin{figure}[H]
\centering
\includegraphics[width=0.88\textwidth]{/home/santiago/repositorios/Nandu_SistemadeAdqusicionEMG/EMG_desarrollo/base_de_datos_electrodos/2026-09-23/O_Prueba3_Sujeto1/plot_paper_combined.png}
\caption{Análisis multimodal de la vocal O (O\_Prueba3\_Sujeto1): Espectrograma de audio con pre-énfasis, oscilograma acústico, activación EMG (Supremo Tricanal) y espectrograma muscular RGB.}
\end{figure}

\subsubsection{Vocal U: U\_Prueba3\_Sujeto1}
\begin{figure}[H]
\centering
\includegraphics[width=0.88\textwidth]{/home/santiago/repositorios/Nandu_SistemadeAdqusicionEMG/EMG_desarrollo/base_de_datos_electrodos/2026-09-23/U_Prueba3_Sujeto1/plot_paper_combined.png}
\caption{Análisis multimodal de la vocal U (U\_Prueba3\_Sujeto1): Espectrograma de audio con pre-énfasis, oscilograma acústico, activación EMG (Supremo Tricanal) y espectrograma muscular RGB.}
\end{figure}

\subsubsection{Vocal E: E\_Prueba3\_Sujeto1}
\begin{figure}[H]
\centering
\includegraphics[width=0.88\textwidth]{/home/santiago/repositorios/Nandu_SistemadeAdqusicionEMG/EMG_desarrollo/base_de_datos_electrodos/2026-09-23/E_Prueba3_Sujeto1/plot_paper_combined.png}
\caption{Análisis multimodal de la vocal E (E\_Prueba3\_Sujeto1): Espectrograma de audio con pre-énfasis, oscilograma acústico, activación EMG (Supremo Tricanal) y espectrograma muscular RGB.}
\end{figure}

\subsection{Prueba4}
\subsubsection{Vocal O: O\_Prueba4\_Sujeto1}
\begin{figure}[H]
\centering
\includegraphics[width=0.88\textwidth]{/home/santiago/repositorios/Nandu_SistemadeAdqusicionEMG/EMG_desarrollo/base_de_datos_electrodos/2026-09-23/O_Prueba4_Sujeto1/plot_paper_combined.png}
\caption{Análisis multimodal de la vocal O (O\_Prueba4\_Sujeto1): Espectrograma de audio con pre-énfasis, oscilograma acústico, activación EMG (Supremo Tricanal) y espectrograma muscular RGB.}
\end{figure}

\begin{figure}[H]
\centering
\begin{subfigure}{0.32\textwidth}
    \includegraphics[width=\textwidth]{/home/santiago/repositorios/Nandu_SistemadeAdqusicionEMG/EMG_desarrollo/base_de_datos_electrodos/2026-09-23/O_Prueba4_Sujeto1/canal_0/avg_lider.png}
    \caption{Canal 0: masetero}
\end{subfigure}
\hfill
\begin{subfigure}{0.32\textwidth}
    \includegraphics[width=\textwidth]{/home/santiago/repositorios/Nandu_SistemadeAdqusicionEMG/EMG_desarrollo/base_de_datos_electrodos/2026-09-23/O_Prueba4_Sujeto1/canal_1/avg_lider.png}
    \caption{Canal 1: -}
\end{subfigure}
\hfill
\begin{subfigure}{0.32\textwidth}
    \includegraphics[width=\textwidth]{/home/santiago/repositorios/Nandu_SistemadeAdqusicionEMG/EMG_desarrollo/base_de_datos_electrodos/2026-09-23/O_Prueba4_Sujeto1/canal_2/avg_lider.png}
    \caption{Canal 2: orbi}
\end{subfigure}
\caption{Forma de onda promedio con su intervalo de dispersión ($\mu \pm \sigma$) y envolvente RMS media para O\_Prueba4\_Sujeto1.}
\end{figure}

\subsubsection{Vocal U: U\_Prueba4\_Sujeto1}
\begin{figure}[H]
\centering
\includegraphics[width=0.88\textwidth]{/home/santiago/repositorios/Nandu_SistemadeAdqusicionEMG/EMG_desarrollo/base_de_datos_electrodos/2026-09-23/U_Prueba4_Sujeto1/plot_paper_combined.png}
\caption{Análisis multimodal de la vocal U (U\_Prueba4\_Sujeto1): Espectrograma de audio con pre-énfasis, oscilograma acústico, activación EMG (Supremo Tricanal) y espectrograma muscular RGB.}
\end{figure}

\begin{figure}[H]
\centering
\begin{subfigure}{0.32\textwidth}
    \includegraphics[width=\textwidth]{/home/santiago/repositorios/Nandu_SistemadeAdqusicionEMG/EMG_desarrollo/base_de_datos_electrodos/2026-09-23/U_Prueba4_Sujeto1/canal_0/avg_lider.png}
    \caption{Canal 0: masetero}
\end{subfigure}
\hfill
\begin{subfigure}{0.32\textwidth}
    \includegraphics[width=\textwidth]{/home/santiago/repositorios/Nandu_SistemadeAdqusicionEMG/EMG_desarrollo/base_de_datos_electrodos/2026-09-23/U_Prueba4_Sujeto1/canal_1/avg_lider.png}
    \caption{Canal 1: -}
\end{subfigure}
\hfill
\begin{subfigure}{0.32\textwidth}
    \includegraphics[width=\textwidth]{/home/santiago/repositorios/Nandu_SistemadeAdqusicionEMG/EMG_desarrollo/base_de_datos_electrodos/2026-09-23/U_Prueba4_Sujeto1/canal_2/avg_lider.png}
    \caption{Canal 2: orbi}
\end{subfigure}
\caption{Forma de onda promedio con su intervalo de dispersión ($\mu \pm \sigma$) y envolvente RMS media para U\_Prueba4\_Sujeto1.}
\end{figure}

\subsection{Prueba5}
\subsubsection{Vocal O: O\_Prueba5\_Sujeto1}
\begin{figure}[H]
\centering
\includegraphics[width=0.88\textwidth]{/home/santiago/repositorios/Nandu_SistemadeAdqusicionEMG/EMG_desarrollo/base_de_datos_electrodos/2026-09-23/O_Prueba5_Sujeto1/plot_paper_combined.png}
\caption{Análisis multimodal de la vocal O (O\_Prueba5\_Sujeto1): Espectrograma de audio con pre-énfasis, oscilograma acústico, activación EMG (Supremo Tricanal) y espectrograma muscular RGB.}
\end{figure}

\begin{figure}[H]
\centering
\begin{subfigure}{0.32\textwidth}
    \includegraphics[width=\textwidth]{/home/santiago/repositorios/Nandu_SistemadeAdqusicionEMG/EMG_desarrollo/base_de_datos_electrodos/2026-09-23/O_Prueba5_Sujeto1/canal_0/avg_lider.png}
    \caption{Canal 0: masetero}
\end{subfigure}
\hfill
\begin{subfigure}{0.32\textwidth}
    \includegraphics[width=\textwidth]{/home/santiago/repositorios/Nandu_SistemadeAdqusicionEMG/EMG_desarrollo/base_de_datos_electrodos/2026-09-23/O_Prueba5_Sujeto1/canal_1/avg_lider.png}
    \caption{Canal 1: -}
\end{subfigure}
\hfill
\begin{subfigure}{0.32\textwidth}
    \includegraphics[width=\textwidth]{/home/santiago/repositorios/Nandu_SistemadeAdqusicionEMG/EMG_desarrollo/base_de_datos_electrodos/2026-09-23/O_Prueba5_Sujeto1/canal_2/avg_lider.png}
    \caption{Canal 2: orbi}
\end{subfigure}
\caption{Forma de onda promedio con su intervalo de dispersión ($\mu \pm \sigma$) y envolvente RMS media para O\_Prueba5\_Sujeto1.}
\end{figure}

\subsubsection{Vocal U: U\_Prueba5\_Sujeto1}
\begin{figure}[H]
\centering
\includegraphics[width=0.88\textwidth]{/home/santiago/repositorios/Nandu_SistemadeAdqusicionEMG/EMG_desarrollo/base_de_datos_electrodos/2026-09-23/U_Prueba5_Sujeto1/plot_paper_combined.png}
\caption{Análisis multimodal de la vocal U (U\_Prueba5\_Sujeto1): Espectrograma de audio con pre-énfasis, oscilograma acústico, activación EMG (Supremo Tricanal) y espectrograma muscular RGB.}
\end{figure}

\begin{figure}[H]
\centering
\begin{subfigure}{0.32\textwidth}
    \includegraphics[width=\textwidth]{/home/santiago/repositorios/Nandu_SistemadeAdqusicionEMG/EMG_desarrollo/base_de_datos_electrodos/2026-09-23/U_Prueba5_Sujeto1/canal_0/avg_lider.png}
    \caption{Canal 0: masetero}
\end{subfigure}
\hfill
\begin{subfigure}{0.32\textwidth}
    \includegraphics[width=\textwidth]{/home/santiago/repositorios/Nandu_SistemadeAdqusicionEMG/EMG_desarrollo/base_de_datos_electrodos/2026-09-23/U_Prueba5_Sujeto1/canal_1/avg_lider.png}
    \caption{Canal 1: -}
\end{subfigure}
\hfill
\begin{subfigure}{0.32\textwidth}
    \includegraphics[width=\textwidth]{/home/santiago/repositorios/Nandu_SistemadeAdqusicionEMG/EMG_desarrollo/base_de_datos_electrodos/2026-09-23/U_Prueba5_Sujeto1/canal_2/avg_lider.png}
    \caption{Canal 2: orbi}
\end{subfigure}
\caption{Forma de onda promedio con su intervalo de dispersión ($\mu \pm \sigma$) y envolvente RMS media para U\_Prueba5\_Sujeto1.}
\end{figure}

\end{document}